% =====================================================================
%  Chapter 26 — The Space Telescopes and What They Do Not Carry
% =====================================================================

\chapter{The Space Telescopes and What They Do Not Carry}
\label{ch:space-telescopes}

A proposal for the Moon is only as good as its comparison with what already
exists. The comparison is more favourable than expected in one respect and
less in another, and the row that decides the argument is not the row anyone
would have guessed.

\section{What is funded}

Two space telescopes carry the mid-infrared channel, and both are specified
in enough detail to be compared with a lunar proposal rather than guessed
at.

NASA's \idxw{NEO Surveyor} will operate a 50-centimetre telescope on a halo orbit\idx{Halo orbit} at
Sun-Earth L1, with 4 to 5.2 and 6 to 10 micrometre channels, surveying 45
to 120 degrees of solar elongation\idx{Elongation} and plus or minus 40 degrees of ecliptic
latitude on a roughly 13-day cadence.

The European Space Agency's \idxw{NEOMIR} concept has a 50-centimetre telescope
passively cooled below 60 kelvin with detectors at 40 kelvin, a 6 to 10
micrometre band, a 30 to 70 degree solar elongation\idx{Elongation} field of regard, and no
atmosphere, no dust and no night. It is at the concept phase A stage, with
an early-2030s date.

Both produce the same observable: an integrated thermal flux, from which a
diameter follows. Neither produces a shape, a spin state or a mass.

\section{The white paper that left the thermal channel to the ground}

The third entry is the one that settles the comparison, and it is a 2026
document rather than a legacy concept.

The white papers for a 3.5-metre segmented-mirror robotic space telescope
describe a baseline that spans 0.2 to 1.5 micrometres and that explicitly
excludes the mid-infrared. The stated architecture is that ``coordinated
ground-based mid-infrared telescopes provide the thermal fluxes required to
infer diameter and albedo, while the space mission supplies contemporaneous
reflected-light measurements and observing geometry''.

That is a division of labour decided by the people designing the most
capable space telescope now under study. The reflected-light channel, the
astrometric channel and the rotational channel are kept; the thermal channel
is handed to the ground.

\begin{table}[!htbp]
\centering
\small
\caption{The platform comparison: what each platform produces, the
population it reaches, and the limit that binds it.}
\label{tab:platforms}
\begin{tabularx}{\textwidth}{@{}P{1.28in}P{1.28in}Y@{}}
\toprule
\rowcolor{panel}
\textbf{Platform} & \textbf{Observable and population} & \textbf{Binding
limit} \\
\midrule
NEO Surveyor\idx{NEO Surveyor}, L1, 0.5 m & integrated thermal flux to a diameter; the whole
discovery stream at a 13-day cadence & diameter only; no shape, no spin \\
\rowcolor{panel}
NEOMIR\idx{NEOMIR}, L1, 0.5 m & as above, at 30 to 70 degrees elongation\idx{Elongation}, early 2030s &
as above; and it is not yet a mission \\
3.5 m segmented-mirror space telescope & resolved photometry, astrometry and
rotation; survey plus rapid response & no mid-infrared channel; relies on
ground mid-infrared \\
\rowcolor{panel}
Ground survey and follow-up & astrometry, lightcurve, photometry; thousands
of objects & unresolved; about 15 per cent with spectra \\
Ground adaptive optics & resolved shape & only at close approach; the wrong
phase angle wastes the window \\
\rowcolor{panel}
Ground bistatic radar & three-dimensional shape, spin, delay-Doppler; more
than fifty objects since 2015 & the echo falls as the fourth power of
distance \\
Lightcurve inversion & spin, pole and a convex shape; about 25,000 models &
the convex hull, and low-order shape only \\
\rowcolor{panel}
Low-Earth-orbit characterisation constellation & visible and near-infrared
characterisation & proposed; constrained by orbit and by brightness \\
Lunar surface array & resolved shape, spin, pole and astrometry; a named
target on demand & needs a kilometre-scale baseline; dust; night \\
\bottomrule
\end{tabularx}
\end{table}

\tabref{tab:platforms} is the manager's comparison table, and every row in
it is a platform that exists, is specified, or has been proposed with
enough detail to compare. The lunar row is the last, and it is the only one
whose binding limit is a construction problem rather than an observing
condition.

\section{The gap that has a shape}

The most capable space telescope concept now under study has chosen the
reflected-light channel for itself and left the thermal flux to the ground.

That is a division of labour with a consequence, and it is worth stating
plainly: it is not that nobody will fly the thermal channel. NEO Surveyor
and NEOMIR will. It is that the \emph{resolved, reflected-light, shape}
channel is not currently served by any space observatory, and that the
ground's contribution to it is smaller than a hundred models across the
entire history of the field.

\begin{table}[!htbp]
\centering
\small
\caption{Where the shape channel currently stands.}
\label{tab:shape-channel}
\begin{tabularx}{\textwidth}{@{}P{1.57in}P{0.99in}Y@{}}
\toprule
\rowcolor{panel}
\textbf{Provider} & \textbf{Models} & \textbf{Status} \\
\midrule
Space telescope, reflected light & none in service & the 3.5 m concept
excludes the mid-infrared and carries no shape programme \\
\rowcolor{panel}
Ground adaptive optics & a handful & close approach, good phase angle \\
Interferometry & fewer than ten, with a hundred projected & terrestrial
baselines of 1,200 km at radio wavelengths \\
\rowcolor{panel}
Fly-by imaging & six & one per multi-billion-dollar mission \\
Lightcurve inversion & about 25,000 & convex hulls \\
\bottomrule
\end{tabularx}
\end{table}

The shape channel is the one place where the portfolio is thinnest, and it
is the channel that the accounting chain runs through. The thermal channel
is well served, by two funded missions. The astrometric channel is
excellently served. The shape channel is served by six spacecraft and
several hundred ground images.

\section{What the space comparison decides and what it does not}

The comparison decides the thermal channel, and it decides it against the
Moon: a 0.5-metre telescope at Sun-Earth L1, already funded for the
thermal flux, at 40 kelvin, with no regolith and no dust, is a better
instrument for that channel than anything the proposal of Chapter 25 could
be built.

The comparison does not decide the shape channel, because no space
observatory is assigned to it.

\begin{keytakeaways}
\item Two funded space missions carry the mid-infrared channel, and the
flagship reflected-light space telescope concept has deliberately left that
channel to the ground.
\item No space observatory currently carries the resolved shape channel, and
the ground's contribution to it is a handful of close-range images and a
projected hundred interferometric models.
\item The comparison therefore decides the thermal channel against the Moon
and leaves the shape channel unassigned.
\item The portfolio is not thin everywhere: it is thin in exactly the place
the accounting chain runs through.
\end{keytakeaways}
